\documentclass{article}
\usepackage[T1]{fontenc}
\usepackage{lmodern}
\usepackage{graphicx}
\usepackage{amsmath,amssymb}
\usepackage[a4paper,margin=2cm]{geometry}
\usepackage[absolute,overlay]{textpos}
\setlength{\TPHorizModule}{1mm}
\setlength{\TPVertModule}{1mm}
\usepackage[indonesian]{babel}
\usepackage{hyperref}
\setlength{\emergencystretch}{2em}
% unit: Halaman 3

\begin{document}

% === Halaman 3 (llm) ===

Contoh:
\begin{enumerate}[(i)]
  \item $23 \pmod{5} = 3$ karena $23$ dibagi $5 = 4 + \text{ sisa } 3$
  \item $-41 \pmod{9} = 4$ karena $|-41| \pmod{9} = 5$, dan $9 - 5 = 4$
  \item $17 \equiv 2 \pmod{3}$ karena $3 \mid (17 - 2)$
  \item $-7 \equiv 15 \pmod{11}$ karena $11 \mid (-7 - 15)$
  \item $23$ dan $40$ relatif prima sebab $\text{PBB}(23, 4) = 1$
  \item $4^{-1} \pmod{9} \equiv 7 \pmod{9}$ karena $4 \cdot 7 \equiv 1 \pmod{9}$
  \item $23^{-1} \pmod{10} = -3 \pmod{10}$ karena $23 \cdot (-3) \equiv 1 \pmod{10}$
\end{enumerate}

Latihan: (a) Hitung $-24 \pmod{11} = ?$
(b) $12^{-1} \pmod{5} \equiv ?$

\end{document}
